% ===== PRÉAMBULE CANONIQUE — à coller VERBATIM en tête de CHAQUE .tex =====
\documentclass[11pt,a4paper]{article}
\usepackage[utf8]{inputenc}\usepackage[T1]{fontenc}\usepackage{lmodern}
\usepackage[french]{babel}
\usepackage{amsmath,amssymb,mathtools}\usepackage{icomma}
\usepackage[version=4]{mhchem}          % \ce{...} : équations & formules
\usepackage{siunitx}\sisetup{locale=FR} % \qty{}{}, \si{}, \ang{}
\usepackage{chemfig}                    % \chemfig{...} : structures développées
\usepackage{xcolor}\usepackage{colortbl}
\usepackage{tikz}\usepackage{pgfplots}\pgfplotsset{compat=1.18}
\usepackage[most]{tcolorbox}\usepackage{enumitem}\usepackage{array,booktabs}
\usepackage[a4paper,margin=2.2cm]{geometry}
\usetikzlibrary{arrows.meta,calc,angles,quotes,positioning,patterns,backgrounds,shapes.geometric}
\definecolor{primary}{RGB}{222,49,99}\definecolor{primarydark}{RGB}{150,22,55}
\definecolor{defcol}{RGB}{0,114,178}\definecolor{methcol}{RGB}{52,92,148}
\definecolor{excol}{RGB}{0,135,90}\definecolor{attncol}{RGB}{200,40,55}
\tcbset{boxbase/.style={enhanced,breakable,boxrule=0.9pt,arc=2.5pt,
  left=3mm,right=3mm,top=2.3mm,bottom=2.3mm,fonttitle=\bfseries,coltitle=white}}
% --- boîtes TUTORIEL (noms EXACTS, ne pas renommer) ---
\newtcolorbox{cle}[1][]{boxbase,colback=defcol!6,colframe=defcol,
  title={\textsc{À retenir}\ifx\relax#1\relax\relax\else\textmd{ : \itshape #1}\fi}}
\newtcolorbox{methode}[1][]{boxbase,colback=methcol!7,colframe=methcol,sharp corners,
  title={\textsc{Méthode}\ifx\relax#1\relax\relax\else\textmd{ --- \itshape #1}\fi}}
\newtcolorbox{exemplebox}[1][]{boxbase,colback=excol!5,colframe=excol,
  title={\textsc{Exemple}\ifx\relax#1\relax\relax\else\textmd{ : \itshape #1}\fi}}
\newtcolorbox{piege}[1][]{boxbase,colback=attncol!6,colframe=attncol,
  title={\textsc{Piège}\ifx\relax#1\relax\relax\else\textmd{ : \itshape #1}\fi}}
\newtcolorbox{remarque}[1][]{boxbase,colback=black!4,colframe=black!45,
  title={\textsc{Remarque}\ifx\relax#1\relax\relax\else\textmd{ : \itshape #1}\fi}}
% --- boîtes EXERCICE / CORRIGÉ (noms EXACTS, auto-comptées) ---
\newtcolorbox[auto counter]{exo}[1][]{boxbase,colback=primary!4,colframe=primary,
  title={\textsc{Exercice}~\thetcbcounter\ifx\relax#1\relax\relax\else\textmd{ --- \itshape #1}\fi}}
\newtcolorbox[auto counter]{corrige}[1][]{boxbase,colback=excol!4,colframe=excol,
  title={\textsc{Corrigé}~\thetcbcounter\ifx\relax#1\relax\relax\else\textmd{ --- \itshape #1}\fi}}
\newcommand{\R}{\mathbb{R}}\newcommand{\N}{\mathbb{N}}\newcommand{\Z}{\mathbb{Z}}\newcommand{\ds}{\displaystyle}
% (un \usepackage{fancyhdr}+en-tête est possible ; il est ignoré côté web)
% =========================================================================
\begin{document}
\sisetup{per-mode=symbol}
\begin{exo}[calculer un rendement]
Calcule le rendement $R=\dfrac{n_{\text{réel}}}{n_{\text{théo}}}\times100$.
\begin{enumerate}[label=\alph*)]
  \item $n_{\text{théo}}=\qty{2}{\mole}$, $n_{\text{réel}}=\qty{1.5}{\mole}$.
  \item $n_{\text{théo}}=\qty{0.4}{\mole}$, $n_{\text{réel}}=\qty{0.3}{\mole}$.
  \item $n_{\text{théo}}=\qty{5}{\mole}$, $n_{\text{réel}}=\qty{4}{\mole}$.
\end{enumerate}
\end{exo}

\begin{exo}[prévoir la quantité réelle]
Calcule la quantité réellement obtenue $n_{\text{réel}}=\dfrac{R}{100}\,n_{\text{théo}}$.
\begin{enumerate}[label=\alph*)]
  \item $n_{\text{théo}}=\qty{2}{\mole}$, $R=\qty{90}{\percent}$.
  \item $n_{\text{théo}}=\qty{0.5}{\mole}$, $R=\qty{60}{\percent}$.
  \item $n_{\text{théo}}=\qty{10}{\mole}$, $R=\qty{50}{\percent}$.
\end{enumerate}
\end{exo}

\begin{exo}[rendement en masse]
Le rendement se calcule aussi avec des masses.
\begin{enumerate}[label=\alph*)]
  \item $m_{\text{théo}}=\qty{40}{\gram}$, $m_{\text{réel}}=\qty{30}{\gram}$ : quel rendement ?
  \item $m_{\text{théo}}=\qty{100}{\gram}$, $m_{\text{réel}}=\qty{85}{\gram}$ : quel rendement ?
  \item $m_{\text{théo}}=\qty{50}{\gram}$, $R=\qty{80}{\percent}$ : quelle masse réelle obtient-on ?
\end{enumerate}
\end{exo}

\begin{exo}[du théorique au rendement]
On brûle \qty{2}{\mole} de \ce{H2} : $\ce{2 H2 + O2 -> 2 H2O}$ (\ce{O2} en excès).
\begin{enumerate}[label=\alph*)]
  \item Quelle quantité \emph{théorique} d'eau prévoit-on ?
  \item On n'en recueille que \qty{1.6}{\mole}. Quel est le rendement ?
\end{enumerate}
\end{exo}

\begin{exo}[rendement et reconversion]
Pour $\ce{N2 + 3 H2 -> 2 NH3}$, on part de \qty{1}{\mole} de \ce{N2} (\ce{H2} en excès). On donne
$M(\ce{NH3})=\qty{17}{\gram\per\mole}$.
\begin{enumerate}[label=\alph*)]
  \item Quelle masse \emph{théorique} d'ammoniac prévoit-on ?
  \item Quelle masse obtient-on réellement si le rendement est de \qty{90}{\percent} ?
\end{enumerate}
\end{exo}

\end{document}
